\ssscp{Lexical taboos require synonyms}{02-03-02-02}
\citet{Simons1982} has written about various kinds of lexical
taboos in Austronesian languages
and pointed out that tabooing assumes
and requires the notion of synonyms -- multiple words
and strategies to talk about the thing you aren’t permitted to
say the normal way.
\citet{GrimesMaryott1994} have written about lexical taboos
and named avoidance speech registers in Austronesian languages,
many of which are associated with hunting and fishing taboos.
\citet{Blust2012b} notes that cuscus are an important animal in the region.
Grimes can verify that from personal experience with multiple
ways of hunting cuscus on \ili{Buru},
trading the dried meat,
and folktales (see \frf{02-21}).

\begin{figure}
	\centering
	\frame{\includegraphics[width=0.375\columnwidth]{02-21.jpg}}
	\fcp{Cuscus meat and broth}{02-21}
\end{figure}

\citet{Grimes1991c,Grimes2000a} gives multiple examples of taboo-generated
synonyms on \ili{Buru} to avoid saying the name of the item being hunted,
or an item considered potentially dangerous.
For example:

 \begin{itemize}
	\item \ili{Buru} \it{uran} ‘shrimp,
				lobster, freshwater crayfish’ (< PMP *qudaŋ ‘shrimp,
				crayfish, lobster’) is more often referred to in the mountainous
				interior as \it{sehe}, which is otherwise a verb ‘reverse,
				back up, withdraw’ associated with the movement of such creatures.
	\item Deer, not native to \ili{Buru} but often hunted,
				are sometimes called \it{menʤaŋan} ‘deer’ (Javanese loan,
				reflecting contact issues), but more often referred to as \it{wadun}
				‘back of neck’, the part of a deer prominently in view when it is feeding head down.
				Loans associated with PWMP *uRsa ‘deer’ (\ili{Malay} \it{rusa}) found
				in many other languages of eastern Indonesia are not known on \ili{Buru}.
	\item The \ili{Rana}
				and Wae Sama dialects of \ili{Buru} refer to mosquitos as \it{seŋet}
				(< PMP *səŋət ‘to sting or bite,
				of an insect’), whereas the \ili{Masarete} dialect has lexicalised
				a descriptive avoidance term \it{inhadat} /ii en-hada-t/ ‘the biting thing’.
	\item Snakes in general on \ili{Buru} are referred to as \it{inewet}
				/ii en-newe-t/ ‘the living thing’,
				and reticulating pythons (\it{Malayopython reticulatus}) are
				called \it{isaleu} /ii saa leu/ ‘the thing that goes before’.
				There are other descriptive terms for specific kinds of snakes,
				with no apparent traces of PMP *hulaR or *nipay.
				Grimes has seen pythons in the \ili{Buru} jungle up to 55cm in circumference,
				and also seen normally fearsome hunter-warriors fleeing them in panic.
				This reflects the perceived danger associated with many lexical taboos.
	\item The elusive civet cat (\it{Viverra tangalunga}
				and \it{Paradoxurus hermaphroditus}, both said to be introduced
				into Maluku from the west,
				see \frf{02-22}), is also feared as a territorial spirit near
				the abandoned mountain village of Ena Biloro.
				Elsewhere it is hunted and eaten.
				The civet cat is called \it{kamyoho} in \ili{Masarete},
				and \it{yoho} in the \ili{Rana} dialect,
				but \it{ŋewaet} in the related \ili{Lisela} language,
				which parses to /ŋee wae-t/ ‘wet nose’ -- an avoidance term.
				Reflexes of PMP *musaŋ ‘civet cat
				and similar’ are not known on \ili{Buru}.\footnote{
						On \ili{Semau} Island
						near Timor, \ili{Helong} uses an avoidance term \it{mat-mitaŋ} ‘black eyes’ to
						refer to the civet cat,
						from the raccoon-like dark fur around the eyes,
						its nocturnal behaviour, and ability to see at night.}
 \end{itemize}

\begin{figure}
	\centering
	\includegraphics[width=0.5\columnwidth]{02-22.jpg}
	\fcp{Asian palm civet}{02-22}
\end{figure}

%\begin{figure}
  %\subfigure[Cuscus meat and broth]{\label{fig:02-21}
    %\includegraphics[width=.45\textwidth]{02-21.jpg}}\hspace{5mm}
  %\subfigure[Asian palm civet]{\label{fig:02-22}
    %\includegraphics[width=.45\textwidth]{02-22.jpg}}
%\caption{Native mammals of Wallacea}
%\end{figure}

These are all examples from \ili{Buru}
and the related \ili{Lisela} language spoken on the same island -- and there are more.
The basic idea of synonyms generated through taboos
and avoidance is repeated in languages throughout the region.
No cognates with *kandoRa or *mans(aə)r have been encountered
in \ili{Buru} or \ili{Lisela} for ‘cuscus’,
although \it{mate} ‘marsupial, cuscus’ has been reported for distantly related \ili{Ambelau}.
And \it{maren} ‘cuscus’ has been reported for recently extinct
\ili{Kayeli}, an \tsc{\ili{Ambon}-Seram} language formerly spoken in eastern \ili{Buru}.
(Note the \ili{Ambelau} and \ili{Kayeli} forms for ‘cuscus’,
and many like them in other languages,
is cognate with *mansər,
which \citet{Blust1982,Blust2012b} glosses ‘bandicoot’, not ‘cuscus’.
Bandicoots are not found as far west as \ili{Ambelau}.) Other competing
terms are also reconstructed for ‘cuscus’,
such as Proto-Celebic *kuse ‘bear cuscus’
and *taŋali ‘dwarf cuscus’,
Proto-\ili{Gorontalo}-\ili{Mongondow} *kusoy ‘cuscus’,
Proto-North Halmahera *kusor ‘cuscus’,
Proto-Nuclear-\ili{Meto} *ʔ-mauka-ʔ ‘cuscus’ \citet{Edwards2021},
and several other reconstructions across Sulawesi,
Halmahera and west Papua.
We return to this issue in \crf{ch:Mar}.
